\section{结论}\label{sec:conclusion}

针对问题一，21种组合的乙醇转化率和 C4 烯烃选择性在各自已测范围内均呈平均正向关系，
但8种组合至少存在一项局部下降或持平。附件2所给一次350~℃实验中，20---273 min内
转化率和收率总体下降，而 C4 烯烃选择性仅在36.72\%---40.32\%之间波动；该结果只描述
本次实验，不能推广为所有组合的长期稳定性。

针对问题二，同一组合的配对变化表明升温通常对应更高的转化率和选择性，但低温段仍有
少数例外。同温排名和扣除温度总体差异后的组合偏差表明，不同组合的观测表现存在差异；
非加和剩余则显示不同组合的升温幅度并不一致。由于各条件没有独立重复，这些结果属于
描述性比较；组合差异与温度差异的相对大小还会随纳入范围改变，不能解释为显著性结论
或某种配方成分的独立因果作用。

针对问题三，在附件实测条件中，一般情景推荐A3在400~℃，收率为44.7281\%；严格低于
350~℃时推荐A2在325~℃，收率为17.2643\%。对11个内部缺测单元的两种插值核查均未
发现会改变低温推荐的候选，但插值结果不属于实验观测，也不证明连续温度上的全局最优。

针对问题四，安排A3在375~℃和425~℃、A2在337.5~℃、A3在400~℃独立重复以及A4在
375~℃共5次实验。其中，A3的两个新温度补充高收率区域两侧的信息，A2的新温度检查
严格低温边界，A3重复检查当前最高观测的局部可复现性，A3与A4在375~℃形成同温比较。
337.5~℃须以设备能够稳定设置为执行前提；
这5次实验用于弥补关键数据缺口，不预设未来结果，也不足以建立完整响应面。
